\section{A fixed finite moment determines the exact maximum}
\label{sec:fixed-moment-reconstruction}

The moment limit can be turned into an explicitly prescribed finite
algebraic selection, without constructing the stationary-value polynomial.
The price is a prohibitive uniform height bound and a very large moment
order. This result is an effective finiteness theorem, rather than a
simple product expression or a practical computation of previously
unknown point counts.

For $n\ge3$, put $r=2n$, $d=n(n-1)=2N$, and define the integers
\begin{equation}\label{eq:uniform-height-constants}
 h_n=\max\{8^N,4n+2\},\qquad
 D_n=rd(d-1)^{r-1},\qquad
 H_n=(16h_nd^r)^{r2^rd^r}.
\end{equation}
These constants are deliberately generous.

\begin{lemma}[An explicit annihilator height bound]
\label{lem:uniform-height-bound}
There is a nonzero $q_n\in\mathbb Z[T]$ satisfying
\[
 q_n(M_n)=0,\qquad \deg q_n\le D_n,\qquad
 \max_j|[T^j]q_n|\le H_n.
\]
The bound is for an annihilating polynomial, without a claim that its
primitive irreducible factor has the same coefficient bound.
\end{lemma}
\begin{proof}
In the $r$ real coordinates $(x_i,y_i)$, the constraint
\[
 S=\left\{(x,y):\sum_i(x_i^2+y_i^2)=n\right\}
\]
is a compact connected sphere defined by one integer polynomial.
Each factor $(x_i-x_j)^2+(y_i-y_j)^2$ has coefficient $\ell^1$ norm $8$.
Consequently $\Delta_n$ has degree $d$ and coefficient height at most
$8^N$. The defining sphere polynomial has degree $2$ and coefficient
height $n\le8^N$.

Apply the deformation and resultant construction of
Jeronimo--Perrucci--Tsigaridas~\cite[Proposition~2.7 and the proof of
Theorem~1.1]{JPT} to the polynomial $-\Delta_n$ on this sphere.
Their construction supplies a nonzero integer polynomial vanishing at
$-M_n$, including when the minimizing set has positive dimension.
There are either $s=0$ or $s=1$ active constraints in that construction;
the degree bound is
\[
 \binom rs d^s(d-1)^{r-s}\le rd(d-1)^{r-1}.
\]
The coefficient estimate explicitly derived in their proof is
\[
 (2^{4-r/2}\widetilde H d^r)^{r2^rd^r},\qquad
 \widetilde H=\max\{H_{\rm input},2r+2m_{\rm input}\}.
\]
Here $H_{\rm input}=8^N$ and $m_{\rm input}=1$, so
$\widetilde H=h_n$. Replacing $2^{4-r/2}$ by the larger integer $16$
gives $H_n$. Finally replace $T$ by $-T$ in the annihilator, which
does not change its coefficient magnitudes.
\end{proof}

For positive integers $D,H$, let $\mathcal A(D,H)$ be the finite set
of all distinct positive real roots of nonzero integer polynomials of
degree at most $D$ and coefficient height at most $H$.
Multiplicity and repetition across different polynomials are ignored.

\begin{lemma}[Uniform separation without a factor-height estimate]
\label{lem:universal-resultant-separation}
Set
\[
 B(D,H)=(2D)!\bigl[(D+1)H2^D\bigr]^{2D},\qquad
 \delta(D,H)=\frac1{B(D,H)+1}.
\]
Any two distinct elements of $\mathcal A(D,H)$ are separated by at
least $\delta(D,H)$.
\end{lemma}
\begin{proof}
Choose integer polynomials $f,g$ of positive degrees at most $D$ and
coefficient heights at most $H$, with $f(\alpha)=g(\beta)=0$.
The polynomial
\[
 R(u)=\operatorname{Res}_t(f(t),g(t+u))\in\mathbb Z[u]
\]
is nonzero: over $\mathbb C$ its roots, with multiplicities, are the
finitely many differences between roots of $g$ and roots of $f$.
Thus $\gamma=\beta-\alpha$ is a nonzero root of $R$.

The coefficient $\ell^1$ norm of $g(t+u)$ is at most
$(D+1)H2^D$. Every polynomial entry of the Sylvester matrix has
$\ell^1$ norm bounded by this number, including the entries from $f$.
The matrix has order at most $2D$. Expanding its determinant and using
submultiplicativity of the coefficient $\ell^1$ norm gives
\[
 \|R\|_1\le(2D)![(D+1)H2^D]^{2D}=B(D,H).
\]
Write $R(u)=u^aQ(u)$ with $Q(0)\ne0$. The integer $Q(0)$ has
absolute value at least one, while every coefficient of $Q$ has
absolute value at most $B=B(D,H)$. If $0<|\gamma|<1/(B+1)$, then
\[
 \left|\sum_{j\ge1}[u^j]Q\,\gamma^j\right|
 \le\frac{B|\gamma|}{1-|\gamma|}<1\le|Q(0)|,
\]
contradicting $Q(\gamma)=0$.
This argument permits repeated roots and requires no factorization.
\end{proof}

Use $B_n=B(D_n,H_n)$ and $\delta_n=1/(B_n+1)$, and put
\begin{equation}\label{eq:explicit-fixed-moment-order}
 U_n=\left(\frac{2n}{n-1}\right)^N,\qquad
 m=n-2,\qquad L_n=m(N+m+1),\qquad
 K_n=\left\lceil4U_nL_n(B_n+1)\right\rceil^2.
\end{equation}
The complete coefficient expansion in
\eqref{eq:coefficient-moment}, evaluated at this prescribed $K_n$,
gives the rational number $A_{n,K_n}$. Let
\[
 E_n=\binom{K_nN+n-2}{n-2}.
\]

\begin{theorem}[A finite coefficient-only algebraic selection]
\label{thm:fixed-moment-reconstruction}
For every $n\ge3$, the exact value $M_n$ is the unique element
$\xi\in\mathcal A(D_n,H_n)$ satisfying
\begin{equation}\label{eq:fixed-moment-root-selector}
 \boxed{\qquad A_{n,K_n}\le\xi^{K_n}\le E_nA_{n,K_n}.\qquad}
\end{equation}
For $n=2$, $M_2=4$ directly.
Thus a single finite moment, whose index is fixed from $n$ before any
geometric calculation, determines the exact maximum.
\end{theorem}
\begin{proof}
The identity
\[
 \sum_{i<j}|z_i-z_j|^2
 =n\sum_i|z_i|^2-\left|\sum_i z_i\right|^2\le n^2
\]
and AM--GM give $M_n\le U_n$.
The moment enclosure~\eqref{eq:moment-enclosure} places $M_n$ in
\[
 I_n=[A_{n,K_n}^{1/K_n},(E_nA_{n,K_n})^{1/K_n}].
\]
This interval has length at most $\delta_n/2$. Indeed, abbreviating
$K=K_n$, the elementary bounds
$E_n\le[K(N+m)]^m$, $\log K\le\sqrt K$ for $K\ge1$, and
$\log(N+m)\le N+m$ give
\[
 \frac{\log E_n}{K}
 \le\frac{m[\log K+\log(N+m)]}{K}
 \le\frac{L_n}{\sqrt K}\le\frac{\delta_n}{4U_n}.
\]
The last quantity is at most $1/2$, so $e^x-1\le2x$ implies
$E_n^{1/K}-1\le\delta_n/(2U_n)$. The lower endpoint of $I_n$
is at most $M_n\le U_n$, whence
\[
 |I_n|=A_{n,K}^{1/K}(E_n^{1/K}-1)\le\delta_n/2.
\]
By Lemma~\ref{lem:uniform-height-bound}, $M_n$ belongs to
$\mathcal A(D_n,H_n)$. By
Lemma~\ref{lem:universal-resultant-separation}, no two distinct
elements of that finite set can lie in $I_n$. This proves existence
and uniqueness in~\eqref{eq:fixed-moment-root-selector}.
\end{proof}

To implement the theorem, enumerate the finitely many bounded integer
coefficient arrays, isolate and deduplicate their positive real roots,
and compute the complete expansion of $V_n^{K_n}$. Real algebraic
arithmetic decides the two inequalities in
\eqref{eq:fixed-moment-root-selector} exactly. Factoring a polynomial
containing the selected root and supplying a rational isolating
interval produces the canonical exact-output representation described
earlier. No enumeration of stationary shapes or construction of their
critical-value polynomial is required by this procedure.

The constants make this procedure prohibitively expensive, even for
small $n$. The sharper degree estimate $n^{2n-4}$ from
Theorem~\ref{thm:exact-algorithm} remains useful; the larger $D_n$ here
is used only to obtain an explicit height bound. The theorem replaces
an unspecified limiting process by finite algebraic selection, and
does not supply a short closed expression or an executed new maximum
for $n\ge8$. The standard-library checker
\texttt{algebraic\_verify.py} verifies small exact moment examples,
translation resultants including repeated roots, and the prescribed
moment-order inequality on manageable constants. It does not execute
the universal polynomial enumeration or the worst-case moment order.

% Bibliography entry to insert in the main document:
% \bibitem{JPT}
% G.~Jeronimo, D.~Perrucci and E.~Tsigaridas,
% \emph{On the minimum of a polynomial function on a basic closed
% semialgebraic set and applications},
% SIAM Journal on Optimization \textbf{23} (2013), no.~1, 241--255.
% \href{https://doi.org/10.1137/110857751}{doi:10.1137/110857751}.
